\section{模型评价与适用范围}
\subsection{方法价值及可推广的处理思路}
四问采用共同的状态定义、径向通量离散与终点判据。有限体积法保留柱坐标权重和有限表面交换阻力，随体坐标将均匀收缩纳入同一求解框架。问题二、三共享初边值解，问题四通过更新物性、半径和输出位置描述移动域过程。

最大含水率判据覆盖径向较慢脱水的位置，同物性固定半径对照则区分总体时间差异与几何影响。解析算例、归一化水分收支以及空间和时间加密，从不同方面检验离散方法及终点稳定性。

\subsection{精度、工况与结论边界}
结论适用于表\ref{tab:closure}给出的理想模型。题目未提供内部温度和含水率观测，现有证据支持方程求解的数值一致性，尚缺少与内部观测比较的预测验证。

长期结果还依赖4小时后的环境延拓。所列情景给出了这一约定附近的变化范围，更广泛工况下的结果需按相应边界重新求解。

本题以含水率0.15为完成目标，计算得到温湿场和达标时间。能耗、药效成分及综合品质未作为目标或约束纳入模型。

\section{结论}

采用守恒有限体积法与隐式积分，得到四问的温湿场。问题一在0.50小时的中心、表面温度分别为33.5753、36.7856摄氏度，含水率分别为2.5500、1.5102，明显失水集中在外层。问题二在3小时的径向温差为0.1169摄氏度，中心、表面含水率仍分别为1.7662、1.0081，水分均匀化滞后于升温。

以径向最大含水率低于0.15为完成条件，问题三、四的越阈时间分别为57.4728、51.0908小时，问题四终点半径为1.2000厘米。全过程结果按规定的时间和空间间隔输出，收缩后的域外位置留空。

问题四相对问题三缩短6.3819小时，降幅为11.1043\%。保持问题四物性而固定半径时，对照时间为129.8449小时，同物性下的收缩降幅为60.6524\%。沿先改变整组物性、再引入收缩的比较路径，时间分别增加72.3721小时、减少78.7541小时，解释了最终时间排序。
